\subsection{Модуль <<Квадрокоптер>>: Pioneer SDK}

Для управления квадрокоптером Pioneer используется специальная библиотека \texttt{pioneer\_sdk}. Она позволяет отправлять команды (взлет, посадка, полет в точку) и получать данные (заряд батареи, видеопоток).

\subsubsection*{Установка библиотеки}
Для начала работы необходимо установить библиотеку через терминал:

\begin{codebox}[Терминал]
\texttt{pip install pioneer\_sdk}
\end{codebox}

\subsubsection*{Подключение к дрону}
Дрон создает Wi-Fi сеть (обычно \texttt{Pioneer\_...}). Подключитесь к ней с вашего компьютера.

\begin{codefile}[python]{python/examples/08_drone_connect.py}
\end{codefile}

\begin{itemize}
    \item \texttt{Pioneer()} — создает объект для управления дроном.
    \item \texttt{drone.get\_battery\_status()} — возвращает напряжение батареи в Вольтах.
    \item \texttt{drone.close\_connection()} — разрывает связь, освобождая порт для других программ.
\end{itemize}

\begin{taskbox}[Диагностика]
Напишите скрипт, который подключается к дрону, проверяет заряд батареи и, если он ниже 3.6 В, выводит предупреждение <<Low Battery!>>.
\end{taskbox}
